\documentclass[english]{article}

\usepackage{ae,aecompl}
\usepackage[T1]{fontenc}
\usepackage[latin9]{inputenc}
\usepackage{textcomp}
\usepackage{graphicx}
\usepackage{eurosym}
\usepackage{hyperref}
\usepackage{float}
\usepackage{fancyhdr}
\pagestyle{fancy}
\lhead{PowerEnJoy RASD}

\newcommand{\carsharing}{\textit {car sharing }}
\newcommand{\powerenjoy}{\textit{PowerEnJoy }}
\newcommand{\registereduser}{\textit {registered user }}
\newcommand{\staff}{\textit{staff }}
\newcommand{\service}{\textit{service }}
\newcommand{\safearea}{\textit{safe area }}
\newcommand{\powergrid}{\textit{Power Grid }}
\newcommand{\resevation}{\textit{reservation }}
\newcommand{\stand}{\textit{stand }}

\makeatletter
\usepackage{babel}


\makeatother

\usepackage{babel}
\begin{document}
\begin{figure}
	\centering
	\includegraphics[scale=0.5]{logo.pdf} 
\end{figure}


\title{PowerEnJoy\\
 Requirement Analysis and Specification Document\\
}

\date{A.A 2016/2017}

\author{Erba Alessandro\\
 Leveni Filippo\\
 Lodi Luca}

\maketitle
\pagebreak{}

\tableofcontents{} \pagebreak{}

\section{Introduction}
	\subsection{Purpose }
	\subsection{Scope and Problem Description}
	\subsection{Goals}
	\subsection{Domain Assumptions }
	\subsection{Glossary}
		\subsubsection{Car Sharing}
			\carsharing is a model of car rental where people rent cars for short periods of time, often by the hour.
		\subsubsection{PowerEnJoy}
			\powerenjoy is the \carsharing brand object of this document. 
		\subsubsection{Registered User}
			A \registereduser is a person subscribed to \powerenjoy which can access to the \carsharing through his smart-phone.
		\subsubsection{Staff}
			\staff is the set of people that are \registereduser but can perform special operations. They are divided in three different categories.
			\begin{itemize}
				\item {Field Staff}
				\begin{itemize}
				\item They own a passepartout for cars that has all kind of issues, for example they can menage cars with low battery.
				\end{itemize}
				\item{Emergency Staff}
				\begin{itemize}
				\item They manage users issues. They are available to communicate with users.
				\end{itemize}
				\item{Management Staff}
				\begin{itemize}
				\item They manage the system, for example they can change system parameters. They can add new \safearea, modify \carsharing prices.
				\end{itemize}
			\end{itemize}
		\subsubsection{Service}
			The \service is the group of actions that a \registereduser can fulfill trough the mobile.	
		\subsubsection {Safe Area}
			A \safearea is a  geographical place\footnote{defined by a set of GPS positions} on the map in which parking is allowed. \safearea are saved into the system. A \registereduser can end a \carsharing parking his car only in safe areas.
	\subsubsection{Power Grid}
		A \powergrid is the station that charge the battery of an electric car. All \powergrid are displayed in the map on-board \powerenjoy cars. Each \powergrid is in a \safearea. A \registereduser who plugs his car into a \powergrid receives the discount expected within \powerenjoy rules.
	\subsubsection{Reservation}
		A \resevation is the possibility that a \registereduser has to book a car at the latest for one hour.
	\subsubsection{Stand}
		A \stand is the possibility that a \registereduser has to pause his \carsharing. During a \stand the \registereduser can leave the car, than he can re-unlock it through his samrt-phone 

\subsection{Further Developments}
	We think that one of the most important things is to release a Service that can be extended and improved. It's trivial that a service like this must be able to be change and be updated.
	For this purpose we will develop a system open to new extensions. \\
	Here we provide a series of possible extensions that could be necessary in the future.
	E.g.:
	\begin{itemize}
		\item New kind of discount policies can be introduced to the system, for example discounts for new users. 
		\item New system policies may be necessary, for example new \carsharing conditions.
		\item New kind of vehicles can be introduced, for example electric scooter and electric bikes.
		\item Support to new payment methods, for example Paypal\textregistered, NFC\footnote{Near Field Communication; NFC devices are used in contactless payment systems, similar to those used in credit cards and electronic ticket smartcards and allow mobile payment to replace/supplement these systems. \href{https://en.wikipedia.org/wiki/Near_field_communication}{NFC on Wikipedia}}.
		\item Develop new services in order to increase the user experience, for example it may be useful to allow communication among users in order to share a car ride and save money.
	\end{itemize}
	\subsection{Used Tools}
	\begin{itemize}
		\item \LaTeX\\
		\item GitHub\\
	\end{itemize}
\section{Specific Requirements}
	\subsection{Functional Requirements}
	\subsection{Non Functional Requirements}

\section{Scenarios Identifying}
\section{UML Models And Use Cases}
	\subsection{Use Cases Diagram}
	\subsection{Actors Identifying}
	\subsection{Use Cases}

	\subsection{Class Diagram}
\section{External Interfaces}
\section{Alloy Model}
\section{Hours of Work}

\end{document}